\documentclass[11pt,a4paper]{article}
\usepackage[T1]{fontenc}
\usepackage{lmodern,amsmath,amssymb,amsthm,microtype}
\usepackage[margin=26mm]{geometry}
\usepackage[colorlinks=true,linkcolor=blue,urlcolor=blue]{hyperref}
\newtheorem{theorem}{Theorem}[section]
\newtheorem{lemma}[theorem]{Lemma}
\newcommand{\Q}{\mathbb Q}\newcommand{\Z}{\mathbb Z}
\title{An intrinsic pole test for elliptic index divisibility}
\author{}\date{17 November 2026}
\begin{document}\maketitle

\begin{abstract}
For a fixed rational congruent-number curve, a reciprocal square
abscissa supplies a balanced pole test at the odd denominator primes
of an elliptic multiple. Under weak rational $abc$ with exponent
$4/3$, every solution outside a fixed finite exceptional set forces
divisibility of the two elliptic indices. Conversely, every nonzero
integer ratio occurs for infinitely many index pairs satisfying the
complete intrinsic-pole preparation. Adding a local logarithmic-ratio
test preserves completeness but does not identify an external
rational input with the index ratio.
\end{abstract}

\section{The source}

Fix $d\in\Z_{>0}$ and
\[
                         E_d:y^2=x^3-d^2x.
\]
Let $P\in4E_d(\Q)$ have infinite order. Choose its sign and replace
it by a fixed positive multiple so that it belongs to the formal group at two with
parameter $\zeta(P)=-x(P)/y(P)$ of valuation at least two.
Write
\[
              H=[h]P,\qquad U=[j]P,\qquad h,j\in\Z\setminus\{0\}.
\]
These index representations are hypotheses of the arithmetic statement.
For a rank-one curve they may be enforced by restricting both points
to a fixed multiple of $E_d(\Q)$. This restriction is a finite
isogeny cover and introduces no further free point coordinate.

Let $q_H\in\Q^\times$ satisfy $q_H^2=x(H)$. There is a positive
choice $q_0=\sqrt{x(P)}\in\Q$. Choose a fixed positive rational
$L$ with
\[
                  v_2(L)=2,\qquad v_7(L)=-1,
\]
and sufficiently large real absolute value, as specified below. Set
$c=q_0L^2$ and define
\begin{equation}\label{eq:weight}
                       u=\frac c{q_H},\qquad w=\frac{u}{(1+u)^2}.
\end{equation}
Use the test $C_1=f(U)=1/x(U)$. Optionally, for an external rational
input $X$, also use
\[
                 C_2=\frac{\zeta(U)}{\zeta(H)}-X.
\]
For either test list put
\begin{equation}\label{eq:detector}
 D_i=\frac{C_i^2}{1+uC_i^2},\qquad
 W=\prod_i\left(1+\frac{D_i^2}{u}\right),\qquad
 b=4wW^2,\qquad \eta=7b.
\end{equation}
Retain all nonzero denominators. Impose the intrinsic equation
\begin{equation}\label{eq:pole}
 Y^2=\eta(\eta+v^2)^2s^4
             -\eta(\eta-v^2)^2(1-\eta)^2,
                \qquad v,s,Y\in\Q.
\end{equation}
The intrinsic pole theorem \cite{Intrinsic} implies from this equation
that $\eta$ has square reduced denominator.

\section{Soundness}

Let $S$ contain $2,7$, the bad reduction primes, and the numerator
and denominator primes of $c$. Write
\[
       x([n]P)=A_n^2/B_n^2,
       \qquad A_n,B_n>0,\quad\gcd(A_n,B_n)=1.
\]
Under weak rational $abc$ with exponent $4/3$, the simple
primitive-prime theorem of \cite{Simple} gives $n_0$ such that
$B_n$ has a primitive prime of valuation one outside $S$ whenever
$n\ge n_0$.

\begin{theorem}\label{thm:sound}
Assume weak rational $abc$ with exponent $4/3$ and
$|h|\ge n_0$. Every rational solution of \eqref{eq:pole}
in the guarded source satisfies $h\mid j$. The assertion remains
valid when any additional rational tests are included in
\eqref{eq:detector}.
\end{theorem}
\begin{proof}
Choose a simple primitive prime $p\notin S$ of $B_{|h|}$.
Since $c$ is a unit there, $v_p(u)=1$.
More generally, when $\nu=v_p(u)$ is odd,
\[
                  v_p(D_i)=\max(2v_p(C_i),-\nu).
\]
The two terms of $1+uC_i^2$ have distinct valuations, which proves
this formula, with the convention $v_p(0)=+\infty$.
Every factor $1+D_i^2/u$ has valuation at most zero: its two terms
have distinct valuations because $2v_p(D_i)-\nu$ is odd.

If $U$ is not formal at $p$, then $v_p(C_1)\le0$, so
$v_p(D_1)\le0$ and $v_p(W)\le-\nu=-1$. Thus
\[
                  v_p(b)=\nu+2v_p(W)<0
\]
is odd. The same is true of $7b$ because $p\ne7$, contradicting
\eqref{eq:pole}. Therefore $U$ is formal at $p$.
The order of $\widetilde P$ modulo $p$ is exactly $|h|$ by
primitivity. Since $[j]\widetilde P=O$, it follows that $h\mid j$.
\end{proof}

The finite restriction $|h|\ge n_0$ requires no free coordinate.
After the cyclic subgroup has been imposed, it can be expressed by
the nonzero guard
\[
                 \prod_{1\le r<n_0}\bigl(x(H)-x([r]P)\bigr)\ne0.
\]
The affine point guards already exclude $h=0$.

\section{Completeness for every integer ratio}

\begin{theorem}\label{thm:complete}
The fixed constant $c$ may be chosen so that the following holds.
For every nonzero integer $m$, there are infinitely many nonzero
integers $h$ with $|h|\ge n_0$ for which, setting $j=mh$ and
choosing $q_H>0$, the input $b$ in \eqref{eq:detector} satisfies
\[
 0<b<1/7,\qquad b\in\Q_2^{\times2},\qquad v_2(b)\ge4,
 \qquad v_7(b)\ge0,
\]
and has square reduced denominator. This also holds with the
optional test $C_2$ when $X=m$. Consequently the intrinsic
theorem supplies rational solutions of \eqref{eq:pole}.
\end{theorem}
\begin{proof}
Choose $L$ as above with $L^2>100$. Take $h$ in a sufficiently
small fixed congruence neighborhood of one so that $H=[h]P$
is close to $P$ at every prime in $S$. Choose the rational
square-root branch $q_H$ close to $q_0$ at these primes. This
branch is coherent: since $P=2R$ for some fixed rational point
$R$, the rational function
\[
        q([h]R)=\frac{x([h]R)^2+d^2}{2y([h]R)}
\]
is a square root of $x([h]P)$; changing the initial sign of $R$
if necessary makes its value $q_0>0$ at $h=1$. Congruences
may also be imposed on $[h]R$.
At the odd primes of $S$, closeness ensures
$v_p(u)=v_p(L^2)$, which is even. At two and seven, refine it
so that $q_0/q_H$ is a local square. Thus
\[
 u\in\Q_2^{\times2},\quad v_2(u)=4,\qquad
 u\in\Q_7^{\times2},\quad v_7(u)=-2.
\]

The formal group at two gives $v_2(\zeta(H))\ge2$ and
$v_2(\zeta(U))\ge v_2(\zeta(H))$, since $U=[m]H$ and
$m\in\Z$. Hence $v_2(C_1)\ge4$. Also the integral formal
multiplication series gives
\[
       \frac{\zeta([m]H)}{\zeta(H)}-m
                      \in\zeta(H)\Z_2,
\]
so $v_2(C_2)\ge2$ when this test is present. Therefore all
factors of $W$ are two-adic units congruent to one modulo two.
It follows that $b$ is a two-adic square and $v_2(b)=6$.

At seven the terms of $1+uC_i^2$ cannot cancel: in the
equal-valuation case $uC_i^2$ has square residue, whereas
$-1$ is not a square modulo seven. Thus the displayed valuation
formula for $D_i$ remains valid even though $v_7(u)=-2$.
It gives $v_7(D_i)\ge2$, so $W\equiv1\pmod7$ and $v_7(b)=2$.

Consider any other prime. If $v_p(u)$ is even, then
$v_p(b)=v_p(4)+v_p(u)-2v_p(1+u)+2v_p(W)$ is even.
If $\nu=v_p(u)<0$ is odd, the balanced valuation formula gives
$v_p(D_i^2/u)\ge-3\nu>0$, hence $W\equiv1\pmod p$ and
$v_p(b)=-\nu>0$.
It remains to consider odd positive $\nu$ outside $S$.
Here $H$ is formal and $v_p(\zeta(H))=\nu$, while $U=[m]H$.
Thus $v_p(C_1)\ge2\nu$ and, if present,
$v_p(C_2)\ge\nu$ by the same integral formal multiplication
series. Every factor of $W$ is a unit congruent to one, and
$v_p(b)=\nu>0$. This proves square denominator at all primes.

Finally take $[h]R$ in a sufficiently small real neighborhood
of $R$. Then $q_H>0$ and $u$ remains larger than $50$.
All tests are real, and
\[
       0\le D_i<1/u,\qquad
       1\le W\le(1+u^{-3})^2.
\]
Consequently $0<b\le4u^{-1}(1+u^{-3})^4<1/7$.
Every fixed arithmetic progression of multiples of a nontorsion
real elliptic point is dense in the appropriate connected real
component. Therefore this real neighborhood is reached for
infinitely many arbitrarily large $h$ in the specified congruence
class. The nonzero affine guards exclude at most finitely many
of these indices. Apply the intrinsic completeness theorem to
$\eta=7b$.
\end{proof}

\section{Scope of the index test}

On a rational elliptic curve each point uses one retained rational
abscissa coordinate and finite algebraic roots. The pole contributes
two free parameters, while its ordinate is a finite quadratic root.
Finite isogeny covers and the displayed nonzero guards introduce no
additional free parameters in the finite-algebra encoder.
An eventual complete polynomial construction must implement this
finite isogeny cover explicitly; no complete polynomial census is
claimed here.

For the square-abscissa chart, retain $q_H,q_U$ freely and use
the monic roots $l_H^2=q_H^4-d^2$, $l_U^2=q_U^4-d^2$, with
$y(H)=q_Hl_H$ and $y(U)=q_Ul_U$. At the prepared rational
values all these point characters are real, while $u=c/q_H>50$
is independent of their signs. Thus every balanced denominator,
$W$, and $b$ is nonzero at every point character. A finite isogeny
cover does not change the images of these point characters.
The intrinsic density theorem allows the pole parameters to avoid
the radicand and division guards simultaneously at every character.

The theorem establishes an integer \emph{index ratio}. With the
optional comparison test, it supplies local congruences between
$X$ and $j/h$. It does not prove $X=j/h$. For fixed $H,U$, the
weight is independent of $X$, and the prepared source persists
under sufficiently small rational perturbations of $X$ at the
real place and finitely many finite places. Thus noninteger
external inputs can still occur. No defining-variable bound for
$\Z\subset\Q$ follows from this source alone.

\begin{thebibliography}{9}
\bibitem{Simple}\emph{Simple primitive denominator primes under weak
$abc$}, companion paper, \nolinkurl{simple_primitive_denominators_v59.tex}.
\bibitem{Intrinsic}\emph{A two-coordinate pole test of coefficient
height four}, companion paper, \nolinkurl{intrinsic_pole_20261117_v56.tex}.
\end{thebibliography}
\end{document}
